\chapter{有限维}\label{ch:b1:findim}

由有限个向量张成的空间具有一个良定义的
\emph{维数} --- 即其所有基共同的大小。下面的证明
均源于同一组合工具，即交换引理：
线性无关族所含向量的个数永远不会超过生成族的。有了维数，
此后处处要用的工具便随之而来：不完全基定理、
族的秩以及格拉斯曼公式。

\section{基的存在性}

\begin{definition}\label{def:b1:findim:def}
称 $E$ 为\emph{有限维}\index{有限维}，当它具有
一个有限生成族时。（否则称为\emph{无限维}：
$K[X]$ 亦属此类，其有限族只能张成
次数有界的多项式。）
\end{definition}

\begin{theorem}[交换引理]\label{thm:b1:findim:exchange}
设 $(g_1, \dots, g_n)$ 生成 $E$，且 $(f_1, \dots, f_p)$ 为
$E$ 中的线性无关族。则 $p \leq n$。
\end{theorem}

\begin{proof}
对 $k \leq p$ 用归纳法证明：\emph{经重新编号
$g_j$ 后，族 $(f_1, \dots, f_k, g_{k+1}, \dots, g_n)$ 生成
$E$} --- 这在每一步都迫使 $k \leq n$，且 $p \leq n$
在最后成立。

$k = 0$ 时即假设。归纳步：设其对 $k - 1$ 成立，则 $f_k$ 是
$(f_1, \dots, f_{k-1},$ $g_k, \dots, g_n)$ 的线性组合。在
此组合中某个 $g_j$（$j \geq k$）的系数非零
---
否则 $f_k$ 将是 $f_1, \dots, f_{k-1}$ 的线性组合，
与线性无关性矛盾。重新编号使 $j = k$，并解出
$g_k$：$g_k$ 是 $(f_1, \dots, f_k, g_{k+1}, \dots,
g_n)$ 的线性组合。$E$ 中每个向量若用 $(k-1)$ 族表示，
便可改用 $k$ 族表示，因此此族生成。
（若 $k - 1 = n$，则已无 $g$ 剩余，$f_k$ 将是
仅为各 $f_i$ 的线性组合，这不可能；故 $k \leq n$。）
\end{proof}

\begin{example}[交换过程的一次演示]
\label{ex:b1:findim:exchangerun}
在 $\R^2$ 中取生成族 $\bigl(g_1, g_2\bigr) =
\bigl((1,0), (0,1)\bigr)$，线性无关族 $\bigl(f_1,
f_2\bigr) = \bigl((1,2), (3,4)\bigr)$。第 $1$ 步：$f_1 = 1\cdot
g_1 + 2\cdot g_2$；$g_2$ 的系数非零，因此
将 $g_2$ 换成 $f_1$：族 $\bigl(f_1, g_1\bigr)$
仍生成（$g_2 = \frac12(f_1 - g_1)$）。第 $2$ 步：$f_2 =
(3,4) = 2\,f_1 + 1\cdot g_1$；剩下的
$g_1$ 的系数非零（必然如此：$f_2$ 不是 $f_1$ 的倍数），
故再次交换：$\bigl(f_1, f_2\bigr)$ 生成 $\R^2$。若
再有第三个线性无关向量 $f_3$，则再无 $g$ 可留下
去吸收它 --- 这正是引理禁止 $3$ 个线性无关
向量位于 $\R^2$ 中的原因。上面的证明即是对此过程的
一般化。
\end{example}

\begin{theorem}[有限维中的基]\label{thm:b1:findim:bases}
设 $E \neq \{0\}$ 为有限维的。
\begin{enumerate}
  \item 从每个有限生成族中都可以抽出一组基。
  \item （不完全基定理\index{不完全基定理}）
        每个线性无关族都可扩充为一组基，所用向量取自任一
        指定的生成族。
  \item $E$ 的所有基都是有限的，且元素个数相同，
        即\emph{维数} $\dim E$。（约定：
        $\dim\{0\} = 0$。）
\end{enumerate}
\end{theorem}

\begin{proof}
(1) 逐个舍去任一可由其余向量线性表出的向量，
族仍然生成，因为在任一用到被舍去向量的表达式中，
都可代入它关于留下向量的线性组合。
过程必然终止 --- 每一步使有限族减少一个，
且恰在剩余向量中再无
可由其余向量线性表出者时停止。最终的族仍然生成，
且线性无关：若有非平凡的零组合，
其中必有某个非零系数，除以它即可解出对应向量，
将其用其余向量表出，从而使该向量
仍可舍去 --- 这与过程已经停止
相矛盾。

(2) 设 $(f_1, \dots, f_p)$ 线性无关，$(g_1, \dots, g_n)$ 生成。
依次考察 $g_1, \dots, g_n$，每次将 $g_j$ 添入当前族，
只要它尚不在当前族的张成之中，
（由\cref{prop:b1:vspaces:freecriteria}~(2)可知族仍线性无关）。
最终的族线性无关且生成：每个 $g_j$ 都落在其
张成之中 --- 要么它已被添入，要么它已是当前族的线性组合。

(3) 两组基各自线性无关且各自生成，由交换引理
可得它们基数之间的两个不等式。（有限性：
基是线性无关的，因此由交换引理，它不大于某个有限
生成族。）
\end{proof}

\begin{example}\label{ex:b1:findim:dims}
$\dim K^n = n$（标准基）；$\dim K_n[X] = n + 1$（单项式）；
$\dim_\R \C = 2$；方程 $y'' + ay' + by = 0$ 的解空间
维数为 $2$（由\cref{thm:b1:diffeq:homogeneous2}可知，解
可由 $(\lambda, \mu) \in K^2$ 双射且线性地参数化）。
\end{example}

\begin{example}[同一集合，两种维数]
\label{ex:b1:findim:twodimensions}
复数对的全体构成的集合 $\C^2$ 是一个 $\C$-向量空间，
维数为 $2$（标准基 $e_1, e_2$） --- 而它同时也是一个
$\R$-向量空间，维数为 $4$，其基为
\[
(1, 0),\quad (\iu, 0),\quad (0, 1),\quad (0, \iu):
\]
任一 $(z, w) = (a + \iu b,\ c + \iu d)$ 都有实坐标
$(a, b, c, d)$，且唯一。维数并非仅由向量集合本身决定：
它计算的是\emph{相对于
允许的标量}的自由度；将标量由 $\C$
缩减为 $\R$ 时，计数便翻一番。（周末问题利用了这种敏感性的
极端情形，即标量一直缩减
到 $\Q$。）
\end{example}

\begin{example}[运行补全算法]
\label{ex:b1:findim:completion}
将线性无关族 $\bigl((1,1,1)\bigr)$ 用标准向量补全为
$\R^3$ 的一组基。运行
\cref{thm:b1:findim:bases}~(2) 的证明，对生成族 $(e_1,
e_2, e_3)$：$e_1 \in \operatorname{Vect}(1,1,1)$ 吗？否（
$(1,1,1)$ 的倍数三个坐标相等），将其加入。$e_2 \in
\operatorname{Vect}\bigl((1,1,1), e_1\bigr)$ 吗？形如
$\alpha(1,1,1) + \beta e_1$ 的组合第二、三个坐标相等，
而 $e_2$ 不然，加入之。族
$\bigl((1,1,1), e_1, e_2\bigr)$ 线性无关，含 $3$ 个向量：停止，
即为一组基（\cref{prop:b1:findim:twoofthree}
将使此做法正式化）。注意结果取决于
$g_j$ 被扫描的先后次序，补全是一个算法，而非一个公式。
\end{example}

\begin{proposition}[三取二规则]\label{prop:b1:findim:twoofthree}
设 $\dim E = n$，$\mathcal{F}$ 为恰含 $n$ 个向量的
$E$ 中的族，则
\[
\mathcal{F} \text{ 线性无关} \iff \mathcal{F} \text{ 生成} \iff
\mathcal{F} \text{ 为基}.
\]
此外，任一线性无关族所含向量个数 $\leq n$，任一生成族
所含向量个数 $\geq n$。
\end{proposition}

\begin{proof}
基数的两个估计来自相对于某组基的替换引理。若
$\mathcal{F}$（大小为 $n$）线性无关而非生成族，存在 $x$ 位于
其张成的子空间之外。添上 $x$ 后得到一个含 $n + 1$
个向量的线性无关族，矛盾。若 $\mathcal{F}$ 是生成族而非线性无关，
抽取一组基（\cref{thm:b1:findim:bases}~(1)），得到含
$< n$ 个向量的基，矛盾。
\end{proof}

\begin{example}[三取二，省去一半的工作]
\label{ex:b1:findim:twoofthreerun}
$\bigl((1,1,0), (0,1,1), (1,0,1)\bigr)$ 是否为 $\R^3$ 的一组基？
计数：三个向量，维数为三，故仅线性无关性
即可判定。零组合给出 $a + c = 0$，$a + b = 0$，$b +
c = 0$；三式相加得 $2(a + b + c) = 0$，再分别从
$a + b + c = 0$ 中减去各原方程，便得 $b = c = a =
0$：线性无关，从而为基；生成性的那一半验证
由定理免费提供。相比之下，
\cref{ex:b1:vspaces:testbasis}，其中同样的双重
验证只能手工完成——一章理论
便换来恰好这一节省，于全书其余
各处每次检验基时皆然。
\end{example}

\begin{method}[计算维数]
\label{met:b1:findim:computedim}
三条标准途径，按使用频率由高到低排列。
\begin{enumerate}
  \item \emph{参数化，再读出一组基。}解出
        定义约束，将一般元素
        用剩下的参数线性表出，并检验
        与各参数相乘的向量线性无关：
        维数即参数个数。（下面在
        $\R^4$ 的一个具体子空间上运行。）
  \item \emph{给出双射的线性参数化。}当
        元素由有限多个值决定时——
        递推关系的初始条件
        （见\cref{exo:b1:findim:10}）、解公式中的系数
        （见\cref{ex:b1:findim:dims}）——维数即
        这些值的个数。
  \item \emph{使用公式。}交与和用 格拉斯曼公式，
        张成用秩，之后求核与像用秩–零化度定理，
        维数通常是
        \emph{计算}出来的，而非猜出来的。
\end{enumerate}
三条途径中，三取二规则都用来收尾：
一旦计数吻合，仅线性无关或仅生成即可下结论。
\end{method}

\begin{example}[路线一，完整演示]
\label{ex:b1:findim:parametrize}
求 $H = \{(x, y, z, t) \in \R^4 : x + y + z + t = 0
\text{ 且 } x = t\}$ 的维数。解：$t = x$，$y + z = -2x$，故 $z =
-2x - y$，其中 $x, y$ 可任取：
\[
(x,\ y,\ -2x - y,\ x) = x\,(1, 0, -2, 1) + y\,(0, 1, -1, 0) .
\]
两个向量线性无关（看前两个坐标：
$(x, y) = (0,0)$），故构成 $H$ 的一组基，且 $\dim H = 2$。
计数本可预料——两个独立的线性约束
在 $\R^4$ 中各应削去一维——但
参数化证明了这一点，\emph{又}给出一组基，
而仅凭预测从来做不到这一点；\cref{exo:b1:findim:9} 使
“每个方程至多削去一维”这一口号化为
定理。
\end{example}

\begin{example}[路线二，完整演示]
\label{ex:b1:findim:route2}
求 $W = \{P \in \R_n[X] : P(1) = P(2) = 0\}$ 的维数（对 $n
\geq 2$）。两次应用因式定理可得
（\cref{thm:b1:poly:factor}；根 $1$ 与 $2$ 互异），
$P \in W$ 当且仅当 $P = (X - 1)(X - 2)\,Q$，其中 $\deg Q \leq
n - 2$。对应 $Q \mapsto (X-1)(X-2)Q$ 是线性的，
可取遍 $W$ 的全部元素，且是单射（乘积为零仅当
$Q = 0$）：$W$ 可由
$\R_{n-2}[X]$ 双射且线性地参数化，于是
\[
\dim W = \dim \R_{n-2}[X] = n - 1 .
\]
参数化自带一组基：取各单项式的像，
即 $\bigl((X-1)(X-2),\ (X-1)(X-2)X,\ \dots,\
(X-1)(X-2)X^{n-2}\bigr)$。每个在
新点处的求值约束恰好削去一维，此即
拉格朗日插值中计数的骨架
（见\cref{thm:b1:poly:lagrange}）。
\end{example}

\begin{example}[积分约束同样削去一维]
\label{ex:b1:findim:integralconstraint}
求 $H = \{P \in \R_2[X] : \int_0^1 P = 0\}$ 的维数。记
$P = a + bX + cX^2$，约束即 $a + \frac b2 + \frac
c3 = 0$；解出 $a$ 并参数化：
\[
P = b\Bigl(X - \frac12\Bigr) + c\Bigl(X^2 - \frac13\Bigr),
\]
于是 $H = \operatorname{Vect}\bigl(X - \frac12,\ X^2 -
\frac13\bigr)$，维数为 $2$（两个多项式次数互异，
故线性无关）。任一线性条件，无论是求值、
积分，还是其他任何线性运算，至多削去一维，
且只要条件不恒为零，便恰好削去一维。
在\cref{ch:b1:linmaps}中将此类运算称为
\emph{线性型}，其解集称为\emph{超平面}；
这里求得的基向量将在
\cref{ch:b1:euclid}中重新出现，作为勒让德族的开端。
\end{example}

\section{子空间、秩与格拉斯曼}

\begin{theorem}[子空间]\label{thm:b1:findim:subspaces}
设 $E$ 为有限维，$F$ 为子空间。则 $F$
为有限维，$\dim F \leq \dim E$，且等号成立当且仅当
$F = E$。此外，每个子空间都存在补子空间。
\end{theorem}

\begin{proof}
$F$ 的线性无关族至多含 $\dim E$ 个向量（由
$E$ 中的交换引理：$F$ 的线性无关族在 $E$ 中亦线性无关，而 $E$
有一个有限生成族）。在 $F$ 的线性无关族中
取一个\emph{极大}者，记其大小为 $p$。这是可行的，因为大小
均为不超过 $\dim E$ 的整数。它生成 $F$：否则存在
$x \in F$ 不属于其张成的子空间，于是添加 $x$ 将得到
$F$ 的大小为 $p + 1$ 的线性无关族，
（由\cref{prop:b1:vspaces:freecriteria}~(2)），与
极大性矛盾。因此它既线性无关又生成，是 $F$ 的基，且
$\dim F = p \leq \dim E$。若 $p = \dim E = n$：$n$ 个
$E$ 中向量构成的线性无关族即为 $E$ 的基（见\cref{prop:b1:findim:twoofthree}），故 $F
\supseteq \operatorname{span} = E$。补空间的构造：
$(f_1, \dots, f_p)$ 为 $F$ 的基，将其扩充为基 $(f_1, \dots, f_p, g_{p+1},
\dots, g_n)$，此即 $E$ 的基（不完全基定理）；令 $G =
\operatorname{Vect}(g_{p+1}, \dots, g_n)$ 满足 $E = F \oplus G$，
（分解的存在性与唯一性，即
大基中的坐标）。
\end{proof}

\begin{definition}[族的秩]\label{def:b1:findim:rank}
称有限向量族的\emph{秩}\index{秩！族的}为
其张成的子空间的维数：
$\operatorname{rk}(x_1, \dots, x_p) = \dim
\operatorname{Vect}(x_1, \dots, x_p) \leq \min(p, \dim E)$,
且秩等于 $p$ 当且仅当该族线性无关。
\end{definition}

\begin{example}[用消元法计算秩]
\label{ex:b1:findim:rankcomputation}
求 $\bigl((1,2,3), (2,3,4), (3,4,5), (1,1,1)\bigr)$ 在
$\R^3$ 中的秩。把某个向量减去其余向量的线性组合，张成空间保持不变
（两族所张成的线性组合集合相同）：
将 $(2,3,4)$ 换为 $(2,3,4) - (1,2,3) = (1,1,1)$，
将 $(3,4,5)$ 换为 $(3,4,5) - (1,2,3) = (2,2,2)$。此时张成空间为
$\operatorname{Vect}\bigl((1,2,3), (1,1,1), (2,2,2),
(1,1,1)\bigr) = \operatorname{Vect}\bigl((1,2,3), (1,1,1)\bigr)$,
这两个向量不成比例，故秩为 $2$。这一
“减去再舍弃”的过程被系统化为
\cref{ch:b1:det}中的高斯消元法。
\end{example}

\begin{example}[通过拼接求和]
\label{ex:b1:findim:sumconcat}
在 $\R^3$ 中取
\[
F = \operatorname{Vect}\bigl((1,2,3),\ (1,1,1)\bigr),
\qquad
G = \operatorname{Vect}\bigl((2,3,4)\bigr) .
\]
和 $F + G$ 由全部
三个生成元拼接而成的族张成，且
\[
(2, 3, 4) = (1, 2, 3) + (1, 1, 1)
\]
表明第三个生成元是多余的：$F + G = F$，维数为 $2$。
等价地 $G \subseteq F$，由上式即可看出。
格拉斯曼公式亦可验证：$\dim(F \cap G) = 2 + 1 - 2 = 1 = \dim G$。
求和的做法是先拼接生成元，
再借助求秩的算法化简该组向量，无需任何新方法。
\end{example}

\begin{theorem}[格拉斯曼公式]\label{thm:b1:findim:grassmann}
设 $F, G$ 为
$E$ 的有限维子空间，则\index{格拉斯曼公式}：
\[
\dim(F + G) = \dim F + \dim G - \dim(F \cap G) .
\]
特别地，$F + G$ 为直和，当且仅当 $\dim(F + G) = \dim F + \dim G$。
\end{theorem}

\begin{proof}
设 $(e_1, \dots, e_r)$ 为 $F \cap G$ 的基，将其
分别扩充为基 $(e_1, \dots, e_r, f_1, \dots, f_s)$（即 $F$ 的基）与
基 $(e_1, \dots, e_r, g_1, \dots, g_t)$（即 $G$ 的基）
（见\cref{thm:b1:findim:bases}~(2)）。断言
\[
\mathcal{B} = (e_1, \dots, e_r, f_1, \dots, f_s, g_1, \dots, g_t)
\]
是 $F + G$ 的基，由计数得 $(r + s) +
(r + t) - r = r + s + t$。

$\mathcal{B}$ 生成 $F + G$：任一 $u + v$（$u \in F$，$v \in G$）
都可由它线性表出。线性无关性：设 $\sum \alpha_i e_i + \sum
\beta_j f_j + \sum \gamma_k g_k = 0$。向量 $w = \sum \gamma_k
g_k = -\sum\alpha_i e_i - \sum\beta_j f_j$ 属于 $G \cap F$，故它
可仅由 $(e_i)$ 线性表出；但 $w$ 也可由 $(g_k)$ 表出，于是在
$G$ 的基下，两种表示必须一致：所有
$\gamma_k = 0$（且 $e$ 的坐标相同）。该关系式化为
$\sum\alpha_i e_i + \sum\beta_j f_j = 0$，这是
$F$ 的基下的一个关系式，于是其余系数全部为零。
\end{proof}

\begin{example}\label{ex:b1:findim:planes}
两个不同的平面 $F, G$（维数 $2$）位于 $\R^3$ 中，满足 $F + G
= \R^3$（它们的和严格包含一个平面），故 $\dim(F \cap G) =
2 + 2 - 3 = 1$：它们总是沿一条直线相交 --- 过原点不存在
“平行平面”。
\end{example}

\begin{example}[格拉斯曼公式的应用，在 $\R^4$ 中]
\label{ex:b1:findim:grassmannrun}
设 $F = \operatorname{Vect}\bigl(e_1,\ e_2,\ (1,1,1,0)\bigr)$，
$G = \operatorname{Vect}(e_3, e_4)$，二者都在 $\R^4$ 中。计算维数：
$\dim F = 3$（第三个生成元的第三坐标非零，
不在 $\operatorname{Vect}(e_1, e_2)$ 中），且 $\dim G = 2$。考虑和：
$F + G$ 包含 $e_1, e_2, e_4$，以及 $e_3 = (1,1,1,0) - e_1 -
e_2$，故其和为整个 $\R^4$。格拉斯曼公式此时 \emph{直接算出}
交的维数，根本无需消元：
\[
\dim(F \cap G) = 3 + 2 - 4 = 1 .
\]
为确定这条直线，在 $G$ 中考察向量 $(0, 0, c, d)$
它属于 $F$ 当且仅当可写成 $a e_1 + b e_2 + \lambda(1,1,1,0)$，
由此推出 $\lambda = -a = -b$ 且 $d = 0$，于是交为
$\operatorname{Vect}(e_3)$ --- 结果一致，因为 $e_3$ 已
在上文中给出为 $F$ 的元素，且按定义属于 $G$。典型的分工是：
格拉斯曼公式预言 \emph{要找多少}，
再由线性方程组确定 \emph{找的是什么}。
\end{example}

\begin{example}[用方程求交]
\label{ex:b1:findim:intersectequations}
若两个子空间均以解集的形式给出，求交只需合并方程。
在 $\R^3$ 中，设 $F = \{x + y + z = 0\}$，且
$G = \{x = y\}$，则
\[
F \cap G = \{x = y,\ 2x + z = 0\}
= \{(x,\ x,\ -2x)\} = \operatorname{Vect}(1, 1, -2),
\]
这是一条直线。用格拉斯曼公式验算：$F + G = \R^3$（两平面互异，
故其和严格包含一个平面），因此
$\dim(F\cap G) = 2 + 2 - 3 = 1$。子空间的两种描述 ---
以方程给出，以生成元给出 --- 各自使一种运算变得平凡：
方程求交即合并方程，生成元求和即拼接生成元；
两种描述之间的互相转换，
正是求解线性方程组的含义（见\cref{ch:b1:det}）。
\end{example}

\begin{remark}[常见误区]
\label{rem:b1:findim:pitfalls}
\emph{维数不沿和相加}，除非和为直和：
$\R^3$ 中两个平面满足 $\dim(F + G) = 3$，而非 $4$；
总应通过交的维数项（格拉斯曼公式）修正。
\emph{包含关系须同时有维数\emph{与}包含本身}：
仅凭 $\dim F = \dim G$ 永远推不出 $F = G$（$\R^2$ 中
两条不同的直线，二者维数相等）。
由\cref{thm:b1:findim:subspaces}的等号情形，须先有 $F \subseteq G$。
\emph{数参数还不是证明}：“$\R^4$ 中的两个方程
，故维数为 $2$”在方程相关时不成立，
例如 $x + y = 0$ 与 $2x + 2y = 0$ 给出维数 $3$；
唯有真正的参数化或秩的计算才能判定。
\emph{不要谈论非子空间的维数}：解集
对\emph{非齐次}方程组的解集不含 $0$；其“维数”
即相应齐次方程组解空间的维数，
（\cref{ch:b1:det}将此精确化）。
\emph{无限维确实存在}：$K[X]$ 中含任意大小的线性无关族
（即单项式），因此不存在
有限生成族 --- 诸如“三取二”规则之类的结论，
只适用于严格有限维的情形，在 $K[X]$ 上会严重失效
（\cref{cor:b1:linmaps:samedim}将对同样的情形给出说明，
即针对单射性与满射性）。
\emph{秩关乎张成的子空间，而非向量组本身}：重复某个向量、
调换次序，或乘以非零常数，秩都不变，
而 $\operatorname{rk} = p$（即向量个数）
是需要\emph{加以证明的性质} --- 它恰好就是线性无关性。一个
$5$ 个向量、秩为 $2$ 的向量族，
含有相当于三个向量的冗余，这一点由
消元法（见\cref{ex:b1:findim:rankcomputation}）明确指出其所在。
\end{remark}

\begin{remark}[维数的用武之地]
\label{rem:b1:findim:whereused}
从此处起，维数是本书最喜爱的计数论证。
\cref{ch:b1:linmaps}证明了秩--零化度定理，
即格拉斯曼公式在线性映射上的版本；\cref{ch:b1:matrices}
用行化简计算秩；\cref{ch:b1:det}把“$n$
个属于 $K^n$ 的向量构成一组基”化为某个数非零。下面的
周末问题展示维数如何做 \emph{算术}：
在域 $\Q$ 上计数维数，可以证明一些无理性命题，
这些命题手算看似无从入手。在本科三年级分册中，
同样的维数计数，经群论细化后，决定哪些
经典尺规作图问题可解 --- 这一故事
正是伽罗瓦理论。
\end{remark}

\begin{remark}[第 3 册内的展望]
\label{rem:b1:findim:perspectives}
维数是本分册其余部分中守恒的量，
值得预先点明其中的守恒定律。
\cref{ch:b1:linmaps}证明了 $\dim E = \dim\ker u +
\operatorname{rk} u$：线性映射压扁的部分
与保留的部分之和，总等于源空间的维数。\cref{ch:b1:det}将此细化
为解集的结构：$p$ 个未知量减去
$\operatorname{rk} A$ 个主元，余下的便是
解空间的维数，高斯消元法将其显示为自由
参数。\cref{ch:b1:euclid}将 $\dim E = \dim F + \dim
F^\perp$ 正交分解，而
\cref{ch:b1:multivar}的周末问题恰好用掉这笔“预算”：$n$ 个数据点，
拟合 $2$ 个参数后，残差占 $n - 2$ 维。
凡在后续章节中计数无法平衡时，
错误必是遗漏了核，或和不是直和 --- 应回头
先检验格拉斯曼公式。
\end{remark}

\section{练习}

\begin{exercise}[$\star$]\label{exo:b1:findim:1}
求下列子空间的一组基及其维数：
\begin{enumerate}
  \item $F = \{(x,y,z) \in \R^3 : x + y + z = 0\}$;
  \item $G = \{(x,y,z,t) \in \R^4 : x = y,\ z = 2t\}$;
  \item $H = \{P \in \R_3[X] : P(1) = P'(1) = 0\}$.
\end{enumerate}
\end{exercise}

\begin{exercise}[$\star$]\label{exo:b1:findim:2}
计算向量族
$\bigl((1,1,1), (1,2,3), (3,5,7), (0,1,2)\bigr)$ 在 $\R^3$ 中的秩，并
抽出其张成的子空间的一组基。
\end{exercise}

\begin{exercise}[$\star$]\label{exo:b1:findim:3}
将线性无关族 $\bigl((1,1,0,0), (0,0,1,1)\bigr)$ 扩充为
$\R^4$ 的一组基，所补向量取标准单位向量，并说明理由。
\end{exercise}

\begin{exercise}[$\star$]\label{exo:b1:findim:4}
证明 $\bigl(1, X, X(X-1), X(X-1)(X-2)\bigr)$ 是
$\R_3[X]$ 的一组基，并求 $X^3$ 在其中的坐标。
\end{exercise}

\begin{exercise}[$\star\star$]\label{exo:b1:findim:5}
设 $F$ 与 $G$ 为维数分别为 $4$ 与 $5$ 的
空间 $E$ 的子空间，且 $\dim E = 7$。$\dim(F \cap
G)$ 的可能取值有哪些？对每个取值，在 $E = \R^7$ 中给出一个实现它的例子。
\end{exercise}

\begin{exercise}[$\star\star$]\label{exo:b1:findim:6}
设 $H = \{P \in \R_n[X] : P(1) = 0\}$。证明 $H$ 是
$\R_n[X]$ 的超平面（即维数为 $n$ 的子空间），并写出
$H$ 的一组基 \emph{（可借助因式定理：$P = (X-1)Q$）}，再
给出一条补直线。
\end{exercise}

\begin{exercise}[$\star\star$]\label{exo:b1:findim:7}
设向量 $u_1, \dots, u_p$ 的秩为 $r$。证明去掉其中一个向量后，
所得向量族的秩为 $r$ 或 $r - 1$，并且
添加一个向量后，秩为 $r$ 或 $r + 1$。由此推得秩
在任意一次增添或删去一个向量时，至多改变 $1$。
\end{exercise}

\begin{exercise}[$\star\star\star$]\label{exo:b1:findim:8}
设 $F_1 \subseteq F_2 \subseteq \dots \subseteq F_k$ 为
$E$ 的子空间（$\dim E = n$），且 $F_i \neq F_{i+1}$（对所有 $i$）。证明 $k
\leq n + 1$。由此推得严格递增的子空间链
$\R^n$ 中此类链的长度至多为 $n + 1$，并给出一条长度最大的
链。
\end{exercise}

\begin{exercise}[$\star\star\star$]\label{exo:b1:findim:9}
设 $E$ 的维数为 $n$，$F$、$G$ 为两个超平面（维数
为 $n - 1$），$F \neq G$。求 $\dim(F \cap G)$。推广之：
$k$ 个超平面之交的维数 $\geq n - k$。
\end{exercise}

\begin{exercise}[$\star\star$]\label{exo:b1:findim:10}
设 $E$ 为满足 $u_{n+2} =
3u_{n+1} - 2u_n$（对所有 $n$）的实数列全体。
\begin{enumerate}
  \item 证明 $E$ 是数列空间的子空间，并且
        $E$ 中的任一数列，完全由数对
        $(u_0, u_1)$ 线性确定；由此推得 $\dim E = 2$。
  \item 验证常数列 $(1)$ 与等比数列
        $(2^n)$ 属于 $E$，且构成 $E$ 的一组基。
  \item 求 $E$ 中满足 $u_0 = 0$，$u_1 = 1$ 的数列。
\end{enumerate}
\end{exercise}

\begin{exercise}[$\star\star$]\label{exo:b1:findim:11}
设 $E$ 的维数为 $n$。
\begin{enumerate}
  \item 若 $F, G$ 为子空间，且 $\dim F + \dim G > n$，证明
        $F \cap G \neq \{0\}$。举例说明：两个子空间，
        维数分别为 $51$ 与 $50$，在 $\R^{100}$ 中总有
        一个非零公共向量。
  \item 若 $H$ 为超平面，$F$ 为子空间，且 $F \cap H =
        \{0\}$，证明 $\dim F \leq 1$。
\end{enumerate}
\end{exercise}

\begin{exercise}[$\star\star\star$]\label{exo:b1:findim:12}
（公共补）设 $F, G$ 为 $E$ 的子空间（有限维），
且 $\dim F = \dim G$。证明 $F$ 与 $G$ 有
\emph{公共补}，即存在子空间 $S$，使 $E = F
\oplus S = G \oplus S$。\emph{（对 $\dim F$ 作向下归纳：若 $F
\neq G \neq E$，取 $x \notin F \cup G$ ---
由\cref{exo:b1:vspaces:12}可知这样的元素存在 --- 并考虑 $F \oplus
\operatorname{Vect}(x)$ 与 $G \oplus \operatorname{Vect}(x)$。）}
\end{exercise}

\section{问题：戴德金塔律}

\begin{problem}\label{pb:b1:findim:1}
第 18--19 章中关于标量的论证，除
域公理（见\cref{def:b1:structures:field}）外均未用到其他性质。因此可取
$K = \Q$，并度量\emph{实数}的集合，
以 $\Q$ 维数为尺度。本题计算
$\Q(\sqrt2, \sqrt3)$ 的维数，证明\emph{戴德金
塔律} $\dim_\Q M = \dim_\Q K \cdot \dim_K M$，并
借纯维数计数得出无理性定理 --- 不用
$\varepsilon$，不用小数，只用基。\index{戴德金塔律}
\index{域！$\R$ 的子域}\index{无理数}

\medskip
\noindent\textbf{第 I 部分 --- 有理标量。}
\begin{enumerate}
  \item 验证 $\Q$ 是域，且第 18--19 章中每个定义与证明
        只用到标量的域公理；由此可知 $\R$ 是 $\Q$ 上的向量空间，
        且交换引理、基定理与格拉斯曼公式在 $\Q$ 上成立。请指出
        唯一一处需用非零标量做除法的步骤，
        它出现在
        \cref{thm:b1:findim:exchange} 的证明中，
        此步即以非零标量作除法。
  \item 证明 $(1, \sqrt2)$ 在 $\Q$ 上线性无关，但在
        $\R$ 上线性相关。设 $\Q(\sqrt2) = \operatorname{Vect}_\Q(1,
        \sqrt2) = \{a + b\sqrt2 : a, b \in \Q\}$，求
        $\dim_\Q \Q(\sqrt2)$。
  \item 证明 $\Q(\sqrt2)$ 对乘法封闭。
        对 $u = a + b\sqrt2$，记 $\sigma(u) = a - b\sqrt2$，
        并记 $N(u) = u\,\sigma(u) = a^2 - 2b^2$。证明 $\sigma(uv) =
        \sigma(u)\sigma(v)$，由此推得 $N(uv) = N(u)N(v)$，并证明
        $N(u) \neq 0$（只要 $u \neq 0$）。
  \item 由此推得每个非零元 $u \in \Q(\sqrt2)$ 都
        在 $\Q(\sqrt2)$ 中有逆元，即 $u^{-1} =
        \sigma(u)/N(u)$：$\Q(\sqrt2)$ 是 $\R$ 的子域。
        计算 $\dfrac1{3 + 2\sqrt2}$ 与 $\dfrac1{1 + \sqrt2}$。
\end{enumerate}

\medskip
\noindent\textbf{第 II 部分 --- 添加 $\sqrt3$。}
\begin{enumerate}[resume]
  \item 证明 $\sqrt6 \notin \Q$ \emph{（比较 $2$ 在
        $6q^2 = p^2$ 两边的指数，
        见\cref{exo:b1:arith:7}）}，再证 $\sqrt3 \notin
        \Q(\sqrt2)$ \emph{（将 $\sqrt3 = a + b\sqrt2$ 平方，分情形讨论
        $ab \neq 0$、$b = 0$、$a = 0$）}。
  \item 记 $M = \operatorname{Vect}_\Q(1, \sqrt2, \sqrt3,
        \sqrt6)$。证明 $M$ 对乘法封闭，
        \emph{（只需列出基向量两两乘积的表格，
        即可）}。
  \item 记 $K = \Q(\sqrt2)$。证明 $(1, \sqrt3)$
        在\emph{域 $K$ 上}线性无关，并由此推得 $M = K +
        K\sqrt3$ 是 $K$-向量空间，维数为 $2$，
        以 $(1, \sqrt3)$ 为基。
  \item 证明 $(1, \sqrt2, \sqrt3, \sqrt6)$ 在
        $\Q$ 上线性无关，从而 $\dim_\Q M = 4$。\emph{（把零关系
        写成 $(a + b\sqrt2) + (c + d\sqrt2)\sqrt3 = 0$，并依次应用
        问题 7 与问题 2。）}
\end{enumerate}

\medskip
\noindent\textbf{第 III 部分 --- 塔律。}设 $\Q \subseteq K
\subseteq M \subseteq \R$，其中 $K$ 与 $M$ 为子域，$(e_1,
\dots, e_m)$ 是 $K$ 作为 $\Q$-向量空间的基，$(f_1,
\dots, f_n)$ 是 $M$ 作为 $K$-向量空间的基。
\begin{enumerate}[resume]
  \item 证明 $mn$ 个乘积 $(e_i f_j)$ 生成 $M$
        （在 $\Q$ 上）。
  \item 证明族 $(e_i f_j)$ 在 $\Q$ 上线性无关。
        \emph{（将 $\Q$ 上的一个零线性组合改写为 $\sum_j
        \bigl(\sum_i \lambda_{ij} e_i\bigr) f_j$ 的形式，其中
        内层系数属于 $K$。）}
  \item 最后给出\emph{戴德金塔律}：
        \[
        \dim_\Q M \;=\; \dim_\Q K \,\cdot\, \dim_K M ,
        \]
        并在 $\Q \subseteq \Q(\sqrt2) \subseteq M$ 上检验此式，
        与问题 7 与问题 8 的结果对照。
  \item （免费得到的域）设 $A \subseteq \R$ 为
        有限维 $\Q$-子空间，且含 $1$，
        对乘法封闭，设 $u \in A$，$u \neq
        0$。证明：若 $(a_1, \dots, a_d)$ 是 $A$ 的基，
        则 $(u a_1, \dots, u a_d)$ 仍是 $A$ 的基；
        由此推得 $u$ \emph{在 $A$ 中}有逆元，即 $A$ 是
        $\R$ 的子域。这一结论恢复了前面
        哪些问题？
  \item 证明：对每个 $x \in A$（沿用问题 12 的记号，$\dim_\Q
        A = d$），族 $(1, x, x^2, \dots, x^{d})$ 线性相关：
        即 $A$ 中每个元素都是某个非零有理系数多项式的根，
        且次数不超过 $d$。
\end{enumerate}

\medskip
\noindent\textbf{第 IV 部分 --- 一个数生成一切。}
设 $s = \sqrt2 + \sqrt3$。
\begin{enumerate}[resume]
  \item 求 $s^2$、$s^3$ 与 $s^4$ 在
        以 $(1, \sqrt2, \sqrt3, \sqrt6)$ 为基的 $M$ 中的坐标。
  \item 证明 $(1, s, s^2, s^3)$ 在 $\Q$ 上线性无关，并由此推得
        $\operatorname{Vect}_\Q(1, s, s^2, s^3) = M$，
        即 $M$ 中每个元素都是 $s$ 的有理系数多项式。试将
        $\sqrt2$ 与 $\sqrt3$ 表示为这样的多项式。
  \item 验证 $s^4 - 10s^2 + 1 = 0$，并证明 $X^4 - 10X^2
        + 1$ 是在 $s$ 处为零的\emph{次数最低的首一多项式}，
        并求出它的四个实根。
  \item 由此推得 $1/s = 10s - s^3$，并指出这个数是什么。证明
        $a + b\sqrt2 + c\sqrt3 + d\sqrt6$（$a,b,c,d \in \Q$）为有理数
        当且仅当 $b = c = d = 0$；特别地，
        $\sqrt2 + \sqrt3 + \sqrt6$ 是无理数
        （是对\cref{exo:b1:reals:7}的加强）。
\end{enumerate}

\medskip
\noindent\textbf{第 V 部分 --- 立方根留在外面。}设 $t
= 2^{1/3}$。
\begin{enumerate}[resume]
  \item 证明 $X^3 - 2$ 没有有理根\emph{（可用
        \cref{exo:b1:poly:5}中的有理根判别法，或用
        赋值计数）}；由此推得 $(1, t)$ 在
        $\Q$ 上线性无关。
  \item 证明 $(1, t, t^2)$ 在 $\Q$ 上线性无关。\emph{（若存在某个非零的
        $P \in \Q[X]$，次数 $\leq 2$，且以 $t$ 为根，取
        其中次数最低者，用它去除 $X^3 - 2$，
        由\cref{thm:b1:poly:division}，推出 $X^3 - 2$
        将有有理根。）}由此推得 $A =
        \operatorname{Vect}_\Q(1, t, t^2)$ 的维数为 $3$，对乘法封闭，
        且为 $\R$ 的子域。
  \item 证明 $t \notin M$：否则 $M$ 将成为
        域 $A$ 上的向量空间，塔律将迫使
        $3 \mid 4$。因此 $2^{1/3}$ 不是
        $1, \sqrt2, \sqrt3, \sqrt6$ 的有理线性组合。
  \item 由此推得 $t \notin \Q(\sqrt2)$，$t + \sqrt2$ 是
        无理数，并且不存在有理数 $a, b$ 满足
        $2^{1/3} = a + b\sqrt3$。
\end{enumerate}

\medskip
\noindent\textbf{第 VI 部分 --- 全部子域与综合。}
\begin{enumerate}[resume]
  \item 证明维数的整除性：若 $\Q \subseteq A
        \subseteq B \subseteq \R$ 为子域，且 $\dim_\Q B$
        有限，则 $\dim_\Q A$ 整除 $\dim_\Q B$。
        $M$ 的子域可能有哪些维数？
  \item 求出 $M$ 的\emph{全部}子域，维数为 $2$
        （在 $\Q$ 上）。\emph{（归结为求 $w = a + b\sqrt2
        + c\sqrt3 + d\sqrt6 \in M \setminus \Q$ 使 $w^2 \in
        \Q$：展开 $w^2$，并令三个无理
        坐标为零。）}
  \item 将 $\dfrac1{1 + \sqrt2 + \sqrt3}$ 表示为以 $(1,
        \sqrt2, \sqrt3, \sqrt6)$ 为基的形式\emph{（乘以适当选取的
        共轭）}。
  \item 综合，用四句话说明：为何 $\Q$ 上的维数是
        实数集合的\emph{算术}不变量；有限维性对
        每个元素迫使什么（问题 13）；
        “三取二”规则如何凭空得出逆元（问题 12）；
        塔律禁止什么（问题 20）。
        指出第 III 部分中所证明的
        定理。
\end{enumerate}
\end{problem}
